% GENERATED FILE. Edit canonical lessons, then run npm run generate:books.
% canonical-source-hash: fnv1a64:72fe927dabeec984
% canonical-lessons: GE-C406-einparken, GE-C406-ausparken, GE-C406-parkluecke, GE-C406-parkhaus, GE-C406-parkschein

\chapter{Einen Platz finden}
\label{ch:ge-einen-platz-finden}
\hlchaptermodality{406}

\begin{chapteropening}
\textbf{By the end of this chapter:} \emph{I can talk about parking: parking, pulling out, the parking space, the car park and the parking ticket.}

Complete the last lesson of chapter 406: I can talk about parking: parking, pulling out, the parking space, the car park and the parking ticket.
\end{chapteropening}

\section[einparken]{einparken — to park (in a space)}

\label{lesson:GE-C406-einparken}



\begin{quote}
Before the new one: say the German for the dent, then the German for to dent.
\end{quote}

\subsection*{You'll want to know: einparken}
\textbf{einparken} — \textquotedblleft{}to park (in a space)\textquotedblright{}.

\begin{quote}
Sie parkt das große Auto ohne Probleme ein.
\end{quote}

\begin{quote}
\emph{She parks the big car without any problems.}
\end{quote}

Separable: \emph{ich parke ein}, \emph{ich habe eingeparkt}. \emph{Parken} alone is more general.

\begin{cousinweb}
\emph{einparken} is built from \emph{ein-} 'in' and \emph{parken} 'to park', which is usually traced to English \emph{to park}. EtymDB does not list the German verb, so the line is hedged.
\end{cousinweb}

\begin{grammarlens}[title={What you've built}]
One more word for telling, explaining and describing at B1.
\end{grammarlens}

\subsection*{Your turn}
\begin{itemize}
\raggedright
  \item \emph{Say it:} \emph{einparken}
  \item \emph{Say it:} \emph{einparken}, once more
  \item \emph{Say it:} a sentence of your own with \emph{einparken}
  \item \emph{Recall:} say the German for the dent, then the German for to dent, then say \emph{einparken} again
\end{itemize}

\begin{tcolorbox}[breakable,colback=teal!4,colframe=teal!35!black,title={Before you move on}]
What does einparken mean? (\textquotedblleft{}To park (in a space)\textquotedblright{}.) Say it once more.
\end{tcolorbox}
\section[ausparken]{ausparken — to pull out (of a parking space)}

\label{lesson:GE-C406-ausparken}



\begin{quote}
Before the new one: say the German for to dent, then the German for to park.
\end{quote}

\subsection*{You'll want to know: ausparken}
\textbf{ausparken} — \textquotedblleft{}to pull out (of a parking space)\textquotedblright{}.

\begin{quote}
Beim Ausparken fährt er fast gegen eine Mauer.
\end{quote}

\begin{quote}
\emph{When pulling out he almost drives into a wall.}
\end{quote}

Separable: \emph{ich parke aus}, \emph{ich habe ausgeparkt}.

\begin{cousinweb}
\emph{ausparken} is built from \emph{aus} 'out' and \emph{parken} 'to park', which is usually traced to English \emph{to park}; EtymDB does not list the German verb.
\end{cousinweb}

\begin{grammarlens}[title={What you've built}]
One more word for telling, explaining and describing at B1.
\end{grammarlens}

\subsection*{Your turn}
\begin{itemize}
\raggedright
  \item \emph{Say it:} \emph{ausparken}
  \item \emph{Say it:} \emph{ausparken}, once more
  \item \emph{Say it:} a sentence of your own with \emph{ausparken}
  \item \emph{Recall:} say the German for to dent, then the German for to park, then say \emph{ausparken} again
\end{itemize}

\begin{tcolorbox}[breakable,colback=teal!4,colframe=teal!35!black,title={Before you move on}]
What does ausparken mean? (\textquotedblleft{}To pull out (of a parking space)\textquotedblright{}.) Say it once more.
\end{tcolorbox}
\section[die Parklücke]{die Parklücke — the parking space (between cars)}

\label{lesson:GE-C406-parkluecke}



\begin{quote}
Before the new one: say the German for to park, then the German for to pull out.
\end{quote}

\subsection*{You'll want to know: die Parklücke}
\textbf{die Parklücke} — \textquotedblleft{}the parking space (between cars)\textquotedblright{}.

\begin{quote}
Wir suchen zwanzig Minuten eine Parklücke.
\end{quote}

\begin{quote}
\emph{We look for a parking space for twenty minutes.}
\end{quote}

Plural: \emph{die Parklücken}. \emph{Die Lücke} — the gap.

\begin{cousinweb}
\emph{Parklücke} is built from \emph{parken} 'to park' and \emph{Lücke} 'gap', inherited from Old High German \emph{luccha} (EtymDB).
\end{cousinweb}

\begin{grammarlens}[title={What you've built}]
One more word for telling, explaining and describing at B1.
\end{grammarlens}

\subsection*{Your turn}
\begin{itemize}
\raggedright
  \item \emph{Say it:} \emph{die Parklücke}
  \item \emph{Say it:} \emph{die Parklücke}, once more
  \item \emph{Say it:} a sentence of your own with \emph{die Parklücke}
  \item \emph{Recall:} say the German for to park, then the German for to pull out, then say \emph{die Parklücke} again
\end{itemize}

\begin{tcolorbox}[breakable,colback=teal!4,colframe=teal!35!black,title={Before you move on}]
What does die Parklücke mean? (\textquotedblleft{}The parking space (between cars)\textquotedblright{}.) Say it once more.
\end{tcolorbox}
\section[das Parkhaus]{das Parkhaus — the multi-storey car park (US: parking garage)}

\label{lesson:GE-C406-parkhaus}



\begin{quote}
Before the new one: say the German for to pull out, then the German for the parking space.
\end{quote}

\subsection*{You'll want to know: das Parkhaus}
\textbf{das Parkhaus} — \textquotedblleft{}the multi-storey car park (US: parking garage)\textquotedblright{}.

\begin{quote}
Das Parkhaus am Bahnhof ist am Wochenende voll.
\end{quote}

\begin{quote}
\emph{The car park at the station is full at the weekend.}
\end{quote}

Plural: \emph{die Parkhäuser}. An underground one is \emph{die Tiefgarage}.

\begin{cousinweb}
\emph{Parkhaus} is built from \emph{parken} 'to park' and \emph{Haus} 'house', inherited from Proto-Germanic \emph{hūsą} (EtymDB).
\end{cousinweb}

\begin{grammarlens}[title={What you've built}]
One more word for telling, explaining and describing at B1.
\end{grammarlens}

\subsection*{Your turn}
\begin{itemize}
\raggedright
  \item \emph{Say it:} \emph{das Parkhaus}
  \item \emph{Say it:} \emph{das Parkhaus}, once more
  \item \emph{Say it:} a sentence of your own with \emph{das Parkhaus}
  \item \emph{Recall:} say the German for to pull out, then the German for the parking space, then say \emph{das Parkhaus} again
\end{itemize}

\begin{tcolorbox}[breakable,colback=teal!4,colframe=teal!35!black,title={Before you move on}]
What does das Parkhaus mean? (\textquotedblleft{}The multi-storey car park (US: parking garage)\textquotedblright{}.) Say it once more.
\end{tcolorbox}
\section[der Parkschein]{der Parkschein — the parking ticket (from a machine)}

\label{lesson:GE-C406-parkschein}



\begin{quote}
Before the new one: say the German for the parking space, then the German for the multi-storey car park.
\end{quote}

\subsection*{You'll want to know: der Parkschein}
\textbf{der Parkschein} — \textquotedblleft{}the parking ticket (from a machine)\textquotedblright{}.

\begin{quote}
Den Parkschein legt er hinter die Windschutzscheibe.
\end{quote}

\begin{quote}
\emph{He puts the parking ticket behind the windscreen.}
\end{quote}

Plural: \emph{die Parkscheine}. It comes from a machine, \emph{der Parkautomat}.

\begin{cousinweb}
\emph{Parkschein} is built from \emph{parken} and \emph{Schein} 'certificate, slip', which first meant 'shine, appearance' and goes back to Proto-Germanic \emph{skīnaną} 'to shine' (EtymDB).
\end{cousinweb}

\begin{grammarlens}[title={What you've built}]
One more word for telling, explaining and describing at B1.
\end{grammarlens}

\subsection*{Your turn}
\begin{itemize}
\raggedright
  \item \emph{Say it:} \emph{der Parkschein}
  \item \emph{Say it:} \emph{der Parkschein}, once more
  \item \emph{Say it:} a sentence of your own with \emph{der Parkschein}
  \item \emph{Recall:} say the German for the parking space, then the German for the multi-storey car park, then say \emph{der Parkschein} again
\end{itemize}

\begin{tcolorbox}[breakable,colback=teal!4,colframe=teal!35!black,title={Before you move on}]
What does der Parkschein mean? (\textquotedblleft{}The parking ticket (from a machine)\textquotedblright{}.) Say it once more.
\end{tcolorbox}
